\documentclass[../../../include/open-logic-section]{subfiles}

\begin{document}

\olfileid[id]{sth}{replacement}{refproofs}
\olsection{Lampiran: Hasil-Hasil seputar Penggantian}

Dalam bagian ini, kita akan membuktikan Refleksi di dalam~$\ZF$. Kita juga
akan membuktikan suatu pengertian yang menjadikan Refleksi ekuivalen dengan
Penggantian. Kemudian kita akan membuktikan suatu konsekuensi menarik dari
semua ini mengenai kekuatan Refleksi/Penggantian. \emph{Peringatan: bagian ini
dengan mudah merupakan bagian matematika paling lanjut dalam buku teks ini.}

Kita mulai dengan sebuah lema yang, demi keringkasan, menggunakan perangkat
notasi berupa \emph{garis atas} untuk menangani barisan variabel atau objek.
Jadi, ``$\overline{a}_k$'' menyingkat ``$a_{k_1}$, \dots, $a_{k_n}$'',
dengan $n$ ditentukan oleh konteks.

\begin{lem}\ollabel{lemreflection}
Untuk setiap $1\leq i\leq k$, misalkan
$\phi_i(\overline{v}_i,x)$ suatu formula. Maka, untuk setiap $\alpha$,
terdapat suatu $\beta>\alpha$ sedemikian sehingga, untuk sebarang
$\overline{a}_1,\ldots,\overline{a}_k\in V_\beta$ dan setiap
$1\leq i\leq k$:
\[
	\exists x\phi_i(\overline{a}_i, x) \rightarrow (\exists x \in V_\beta) \phi_i(\overline{a}_i, x)
\]
\end{lem}

\begin{proof}
Kita mendefinisikan suatu suku~$\mu$ sebagai berikut:
$\mu(\overline{a}_1,\ldots,\overline{a}_k)$ merupakan tahap paling awal~$V$
yang memenuhi semua implikasi berikut, untuk $1\leq i\leq k$:
\[
\exists x\phi_i(\overline{a}_i, x) \rightarrow (\exists x \in V) \phi_i(\overline{a}_i, x)
\]
Mudah dipastikan bahwa $\mu(\overline{a}_1,\ldots,\overline{a}_k)$ ada bagi
semua $\overline{a}_1,\ldots,\overline{a}_k$. Sekarang, dengan menggunakan
Penggantian dan teorema rekursi kita, definisikan:
\begin{align*}
	S_0 & = V_{\alpha+1}\\
	S_{n+1} & = S_n \cup \bigcup
	\Setabs{\mu(\overline{a}_1, \ldots, \overline{a}_k)}
	{\overline{a}_1, \ldots, \overline{a}_k \in S_n} \\
	S &= \bigcup_{m < \omega} S_m.
\end{align*}
Setiap $S_n$, dan karena itu $S$ sendiri, merupakan suatu tahap sesudah
$V_\alpha$. Sekarang tetapkan
$\overline{a}_1,\dots,\overline{a}_k\in S$; maka terdapat suatu
$n<\omega$ sedemikian sehingga
$\overline{a}_1,\dots,\overline{a}_k\in S_n$. Tetapkan suatu
$1\leq i\leq k$, dan andaikan
$\exists x\phi_i(\overline{a}_i,x)$. Berdasarkan konstruksi,
$(\exists x\in\mu(\overline{a}_1,\ldots,\overline{a}_k))
\phi_i(\overline{a}_i,x)$; jadi,
$(\exists x\in S_{n+1})\phi_i(\overline{a}_i,x)$ dan karena itu
$(\exists x\in S)\phi_i(\overline{a}_i,x)$. Dengan demikian, $S$ adalah
$V_\beta$ yang kita cari.
\end{proof}
\noindent
Sekarang kita dapat membuktikan
\olref[sth][replacement][ref]{reflectionschema} secara cukup langsung:

\begin{proof}[Bukti]
Tetapkan $\alpha$. Tanpa mengurangi keumuman, kita dapat mengandaikan bahwa
satu-satunya penghubung dalam~$\phi$ ialah $\exists$, $\lnot$, dan~$\land$
(karena ketiganya cukup secara ekspresif). Misalkan
$\psi_1,\ldots,\psi_k$ mengenumerasi setiap subformula~$\phi$ menurut
kompleksitas, sehingga $\psi_k=\phi$. Berdasarkan \olref{lemreflection},
terdapat $\beta>\alpha$ sedemikian sehingga, untuk sebarang
$\overline{a}_i\in V_\beta$ dan setiap $1\leq i\leq k$:
\begin{align}\label{reflectionnicelybehaved}
	\exists x\psi_i(\overline{a}_i, x) \rightarrow 
	(\exists x \in V_\beta) \psi_i(\overline{a}_i, x)\tag{*}
\end{align}
Dengan induksi pada kompleksitas~$\psi_i$, kita menunjukkan bahwa
$\psi_i(\overline{a}_i)\leftrightarrow
\psi_i^{V_\beta}(\overline{a}_i)$ bagi sebarang
$\overline{a}_i\in V_\beta$. Jika $\psi_i$ atomik, klaim ini jelas.
Bikondisional tersebut juga menetapkan bahwa, jika $\psi_i$ merupakan negasi
atau konjungsi subformula-subformula yang memenuhi sifat ini, maka $\psi_i$
sendiri memenuhi sifat tersebut. Jadi, satu-satunya kasus yang menarik
menyangkut kuantifikasi. Tetapkan $\overline{a}_i\in V_\beta$; maka:
\begin{align*}
	(\exists x \psi_i(\overline{a}_i, x))^{V_\beta}
	&\text{ jika dan hanya jika }
	(\exists x \in V_\beta)\psi_i^{V_\beta}(\overline{a}_i, x)
	&&\text{berdasarkan definisi}\\
	&\text{ jika dan hanya jika }
	(\exists x \in V_\beta)\psi_i(\overline{a}_i, x)
	&&\text{berdasarkan hipotesis}\\
	&\text{ jika dan hanya jika }
	\exists x \psi_i(\overline{a}_i, x)
	&&\text{berdasarkan \eqref{reflectionnicelybehaved}.}
\end{align*}
Ini menyelesaikan induksi; hasilnya mengikuti karena $\psi_k=\phi$.
\end{proof}

Kita telah membuktikan Refleksi di dalam~$\ZF$. Bukti kita pada dasarnya
mengikuti \citet{Montague1961}. Sekarang kita hendak membuktikan di
dalam~$\Z$ bahwa Refleksi mengakibatkan Penggantian. Bukti mengikuti
\citet{Levy1960}, tetapi dengan suatu penyederhanaan.

Karena kita bekerja di dalam~$\Z$, kita tidak dapat menyajikan Refleksi tepat
dalam bentuk di atas. Bagaimanapun, kita merumuskan Refleksi dengan notasi
``$V_\alpha$'', dan notasi itu tidak dapat didefinisikan di dalam~$\Z$
(lihat \olref[sth][spine][zf]{sec}). Karena itu, sebagai gantinya kita
menawarkan suatu rumusan Refleksi yang tampak lebih lemah:

\begin{defish}
\emph{Refleksi-Lemah.} Untuk sebarang formula~$\phi$, terdapat suatu himpunan
transitif~$S$ sedemikian sehingga $0$, $1$, dan setiap parameter~$\phi$
merupakan !!{element}s dari~$S$, serta
$(\forall\overline{x}\in S)(\phi\liff\phi^S)$.
\end{defish}

Untuk menggunakan prinsip ini dalam membuktikan Penggantian, mula-mula kita
mengikuti \citet[bagian pertama Teorema 2]{Levy1960} dan menunjukkan bahwa
kita dapat ``merefleksikan'' dua formula sekaligus:

\begin{lem}[dalam $\Z+\text{Refleksi-Lemah}$.]\ollabel{lem:reflect}
Untuk sebarang formula $\psi,\chi$, terdapat suatu himpunan transitif~$S$
sedemikian sehingga $0$ dan $1$ (serta setiap parameter formula-formula itu)
merupakan !!{element}s dari~$S$, dan
$(\forall\overline{x}\in S)((\psi\liff\psi^S)\land
(\chi\liff\chi^S))$.
\end{lem}

\begin{proof}
Misalkan $\phi$ formula $(z=0\land\psi)\lor(z=1\land\chi)$.

Di sini kita menggunakan singkatan; ``$z=0$'' seharusnya ditulis lengkap
sebagai ``$\forall t\,t\notin z$'', dan ``$z=1$'' sebagai
``$\forall s(s\in z\liff\forall t\,t\notin s)$''. Namun, karena
$0,1\in S$ dan $S$ transitif, formula-formula ini \emph{absolut} bagi~$S$;
artinya, keduanya berlaku pada objek yang sama, baik kuantornya dibatasi
pada~$S$ maupun tidak.\footnote{Secara lebih formal, dengan $\xi$ salah satu
dari kedua formula ini, $\xi(z)\liff\xi^S(z)$.}

Berdasarkan Refleksi-Lemah, kita mempunyai suatu $S$ yang sesuai sedemikian
sehingga:
\begin{align*}
	(\forall z, \overline{x} \in S)(&\phi \liff \phi^S)\\
	\text{yakni }(\forall z, \overline{x} \in S)(&((z = 0 \land \psi) \lor (z = 1 \land \chi)) \liff {}\\
	&\phantom{(}((z = 0 \land \psi) \lor (z = 1 \land \chi))^S)\\
	\text{yakni }(\forall z, \overline{x} \in S)(&((z = 0 \land \psi) \lor (z = 1 \land \chi))\liff {}\\
	&\phantom{(}((z = 0 \land \psi^S) \lor (z = 1 \land \chi^S)))\\
	\text{yakni }(\forall \overline{x} \in S)(&(\psi \liff \psi^S) \land (\chi \liff \chi^S)).
\end{align*}
Klaim kedua mengakibatkan klaim ketiga karena ``$z=0$'' dan ``$z=1$''
absolut bagi~$S$; klaim keempat mengikuti karena $0\neq1$.
\end{proof}\noindent
Sekarang kita dapat memperoleh Penggantian hanya dengan mengikuti dan
menyederhanakan \citet[Teorema 6]{Levy1960}:

\begin{thm}[dalam $\Z$ + Refleksi-Lemah]\label{thm:replacement}
Untuk sebarang formula $\phi(v,w)$ dan sebarang~$A$, jika
$(\forall x\in A)\lexists![y][\phi(x,y)]$, maka
$\Setabs{y}{(\exists x\in A)\phi(x,y)}$ ada.
\end{thm}

\begin{proof}
Tetapkan $A$ sedemikian sehingga
$(\forall x\in A)\lexists![y][\phi(x,y)]$, dan definisikan formula-formula:
\begin{align*}
	\psi &\text{ sebagai } (\phi(x, z) \land A = A)\\
	\chi &\text{ sebagai } \lexists[y][\phi(x, y)].
\end{align*}
Dengan menggunakan \olref{lem:reflect}, karena $A$ merupakan parameter
bagi~$\psi$, terdapat suatu himpunan transitif~$S$ sedemikian sehingga
$0,1,A\in S$ (bersama setiap parameter lain), dan sedemikian sehingga:
\[
	(\forall x,z \in S)((\psi \liff \psi^S) \land (\chi \liff \chi^S))
\]
Jadi, secara khusus:
\begin{align*}
	(\forall  x, z \in S)(&\phi(x, z) \liff \phi^S(x, z))\\
	(\forall x \in S)(&\exists y\phi(x, y) \liff (\exists y \in S)\phi^S(x, y)).
\end{align*}
Dengan menggabungkan keduanya dan mengamati bahwa $A\subseteq S$ karena
$A\in S$ dan $S$ transitif:
\begin{align*}
	(\forall x \in A)(&\exists y\phi(x, y) \liff (\exists y \in S)\phi(x, y)).
\end{align*}
Sekarang $(\forall x\in A)(\lexists![y\in S])\phi(x,y)$, karena
$(\forall x\in A)\lexists![y][\phi(x,y)]$. Berdasarkan Pemisahan,
$\Setabs{y\in S}{(\exists x\in A)\phi(x,y)}
=\Setabs{y}{(\exists x\in A)\phi(x,y)}$.
\end{proof}

\end{document}
